%% 2-8 2回かけると□になる数：平方根（第11回で書き足し。ページ数は自由・偶数）
%% この章で定義した部品（頭文字 bkiii）：
%%   \bkiiinumline{(左はし)}{長さ mm}{左の数}{右の数}{目もりの数}   数直線の1本（はさみうちの図）
\makeatletter
\newcommand\bkiiinumline[5]{\begin{scope}[shift={#1}]
  \draw[bkInk, line width=0.8pt] (0,0) -- (#2,0);
  \foreach \k in {0,...,#5} \draw[bkInk, line width=0.5pt] ({#2*\k/#5},-1.1) -- ({#2*\k/#5},1.1);
  \node[anchor=north, font=\sffamily\footnotesize, inner sep=0pt] at (0,-2.2) {#3};
  \node[anchor=north, font=\sffamily\footnotesize, inner sep=0pt] at (#2,-2.2) {#4};\end{scope}}
% 小さな見出し  \bkiiihead{見出し}
\newcommand\bkiiihead[1]{\par\addvspace{5pt}\noindent{\sffamily\bfseries\color{bkTheme!85!black}▶\ #1}\par\nopagebreak\vspace{1pt}}
\makeatother

\bkchapter{2}{8}{2\R{回}{かい}かけると□になる\R{数}{かず}：\\\R{平方根}{へい|ほう|こん}}{\bklevel{2}}{2\R{回}{かい}かけると□になる\R{数}{かず}：\R{平方根}{へい|ほう|こん}}
\bkmichishirube{1-21,2-2,2-6,2-7}{2-9,2-15,3-10,4-4,6-6}

\begin{goal}
\begin{itemize}[label={\color{bkTheme}●}, leftmargin=1.6\zw, itemsep=0pt, topsep=1pt]
\item \bklevel{2}\enspace \R{正方形}{せい|ほう|けい}の\R{面積}{めん|せき}から1\R{辺}{ぺん}を、2\R{乗}{じょう}ではさんで\R{求}{もと}められる。
\item \bklevel{2}\enspace $\sqrt{\phantom{2}}$ の\R{意味}{い|み}がわかり、\R{電卓}{でん|たく}や\R{計算}{けい|さん}で\R{平方根}{へい|ほう|こん}を\R{求}{もと}められる。
\item \bklevel{2}\enspace \R{振}{ふ}り\R{子}{こ}の\R{周期}{しゅう|き}が、\R{糸}{いと}の\R{長}{なが}さの $\sqrt{\phantom{2}}$ に\R{比例}{ひ|れい}することを\R{説明}{せつ|めい}できる。
\end{itemize}
\end{goal}

\begin{hatena}
\noindent\begin{minipage}[c]{0.43\linewidth}
\R{面積}{めん|せき}が\qty{2}{m^2}の\R{正方形}{せい|ほう|けい}の\R{花}{か}だんを\R{作}{つく}ります。1\R{辺}{ぺん}を\R{何}{なん}mにすればよいでしょう。

1\R{辺}{ぺん}\qty{1}{m}では\qty{1}{m^2}、1\R{辺}{ぺん}\qty{2}{m}では\qty{4}{m^2}。1\R{辺}{ぺん}を\qty{1.5}{m}にすると……？
\end{minipage}\hfill
\begin{minipage}[c]{0.55\linewidth}\centering
\begin{tikzpicture}[x=1.1mm, y=1.1mm]
  \filldraw[fill=bkGreen!25, draw=bkGreen!70!black, line width=0.8pt] (0,0) rectangle (11,11);
  \filldraw[fill=bkTheme!25, draw=bkTheme!80!black, line width=0.8pt] (17,0) rectangle (32.6,15.6);
  \filldraw[fill=bkGreen!25, draw=bkGreen!70!black, line width=0.8pt] (39,0) rectangle (61,22);
  \node[font=\sffamily\footnotesize] at (5.5,5.5) {\qty{1}{m^2}};
  \node[font=\sffamily\footnotesize] at (24.8,7.8) {\qty{2}{m^2}};
  \node[font=\sffamily\footnotesize] at (50,11) {\qty{4}{m^2}};
  \node[anchor=north, font=\sffamily\footnotesize] at (5.5,-1.5) {\qty{1}{m}};
  \node[anchor=north, font=\sffamily\bfseries\small, text=bkTheme!80!black] at (24.8,-1.2) {？\,m};
  \node[anchor=north, font=\sffamily\footnotesize] at (50,-1.5) {\qty{2}{m}};
\end{tikzpicture}
\end{minipage}
\end{hatena}

2\R{乗}{じょう}すると\R{面積}{めん|せき}になる\R{数}{かず}を、\R{逆}{ぎゃく}にさがします。この\R{章}{しょう}では、そんな\R{数}{かず}に\R{名前}{な|まえ}と\R{記号}{き|ごう}をつけ、\R{振}{ふ}り\R{子}{こ}のゆれる\R{時間}{じ|かん}にもかくれていることを\R{見}{み}つけます。

\clearpage
%% ---------------- やってみよう・思い出そう ----------------
\begin{yattemiyou}[\R{振}{ふ}り\R{子}{こ}の\R{周期}{しゅう|き}をはかろう]
\noindent\begin{minipage}[c]{0.66\linewidth}
\begin{enumerate}[label={\bkcircled[bkGreen]{\arabic*}}, leftmargin=2.4\zw, itemsep=1pt, topsep=0pt]
\item \R{糸}{いと}の\R{先}{さき}におもりをつけて、\R{支点}{し|てん}からおもりのまん\R{中}{なか}までを\qty{10}{cm}にする。
\item \R{小}{ちい}さくふらして、10\R{往復}{おう|ふく}する\R{時間}{じ|かん}をストップウォッチではかる。3\R{回}{かい}はかって\R{平均}{へい|きん}しよう（2-6）。
\item 10でわって、1\R{往復}{おう|ふく}の\R{時間}{じ|かん}を\R{求}{もと}める。これを\textbf{\R{周期}{しゅう|き}}という。
\item \R{糸}{いと}の\R{長}{なが}さを4\R{倍}{ばい}の\qty{40}{cm}、9\R{倍}{ばい}の\qty{90}{cm}にして、\R{同}{おな}じようにはかろう。
\end{enumerate}
\end{minipage}\hfill
\begin{minipage}[c]{0.31\linewidth}\centering
\begin{tikzpicture}[x=1mm, y=1mm]
  \fill[bkGray!45] (-9,0) rectangle (9,2.2);
  \fill[bkInk] (0,0) circle (0.6);
  \draw[bkInk!25, line width=0.6pt, dashed] (-90+18:30) arc (-90+18:-90-18:30);
  \shade[ball color=bkOrange!25] (-90+18:30) circle (2.6);
  \shade[ball color=bkOrange!25] (-90-18:30) circle (2.6);
  \draw[bkInk, line width=0.7pt] (0,0) -- (0,-30);
  \shade[ball color=bkOrange] (0,-30) circle (2.6);
  \draw[bkGray, line width=0.5pt, <->] (-4,-0.3) -- (-4,-30);
  \node[anchor=east, font=\sffamily\footnotesize, inner sep=0pt] at (-5.5,-13) {\R{長}{なが}さ};
  \draw[bkOrange!80!black, line width=0.6pt, <->] (-90-14:35.5) arc (-90-14:-90+14:35.5);
  \node[anchor=north, font=\sffamily\footnotesize] at (0,-37.5) {\R{行}{い}って\R{帰}{かえ}って1\R{往復}{おう|ふく}};
\end{tikzpicture}
\end{minipage}
\par\vspace{4pt}
\begin{bkcenter}
\begin{bktable}{colspec={Q[l,m,wd=36mm] *{3}{Q[c,m,wd=20mm]}}, rowsep=2pt, colsep=4pt, row{1}={bg=bkGreen!14}, column{1}={font=\sffamily\small}}
\R{糸}{いと}の\R{長}{なが}さ & \qty{10}{cm} & \qty{40}{cm} & \qty{90}{cm} \\
10\R{往復}{おう|ふく}の\R{時間}{じ|かん}（\R{秒}{びょう}） & 6.3 & 12.7 & 19.0 \\
\R{周期}{しゅう|き}（\R{秒}{びょう}） & 0.63 & 1.27 & 1.90 \\
\end{bktable}
\bktabcaption{はかった\R{例}{れい}。\R{長}{なが}さを4\R{倍}{ばい}・9\R{倍}{ばい}にすると、\R{周期}{しゅう|き}は\R{何倍}{なん|ばい}になったかな？}
\end{bkcenter}
\end{yattemiyou}

\begin{omoidasou}
\bkrecall{1-21}{\R{正方形}{せい|ほう|けい}の\R{面積}{めん|せき}＝1\R{辺}{ぺん}×1\R{辺}{ぺん}}
\bkrecall{2-7}{$3\times3=3^2$（3の2\R{乗}{じょう}）。2\R{乗}{じょう}を\R{平方}{へい|ほう}ともいう。}
\bkrecall{2-2}{マイナス×マイナスはプラス。$(-3)\times(-3)=9$}
\bkrecall{2-6}{\R{四捨五入}{し|しゃ|ご|にゅう}で\R{小数}{しょう|すう}\R{第}{だい}2\R{位}{い}までにすると、1.414は1.41}
\end{omoidasou}

\clearpage
%% ---------------- 数学の道具 ----------------
\bksection{doguu}{\R{数学}{すう|がく}の\R{道具}{どう|ぐ}}

\bksubsection{1}{\R{面積}{めん|せき}から1\R{辺}{ぺん}を\R{求}{もと}める}

\leavevmode\label{bk:2-8:men}1\R{辺}{ぺん}を2\R{乗}{じょう}すると、\R{正方形}{せい|ほう|けい}の\R{面積}{めん|せき}になります。\R{表}{ひょう}を\R{逆}{ぎゃく}に\R{読}{よ}めば、\R{面積}{めん|せき}から1\R{辺}{ぺん}がわかります。

\begin{bkfigure}
\begin{tikzpicture}[x=1mm, y=1mm]
  \foreach \a/\b [count=\i from 0] in {1/1,2/4,3/9,4/16,5/25} {
    \node[bk node, minimum width=13mm, minimum height=7mm, inner ysep=1mm] (t\i) at (24+\i*18,7) {\a};
    \node[bk node, draw=bkBlue, minimum width=13mm, minimum height=7mm, inner ysep=1mm] (m\i) at (24+\i*18,-7) {\b};
    \draw[bk link, bkTheme!80!black, shorten <=1mm, shorten >=1mm] ([xshift=-2.2mm]t\i.south) -- ([xshift=-2.2mm]m\i.north);
    \draw[bk link, bkBlue!80!black, shorten <=1mm, shorten >=1mm] ([xshift=2.2mm]m\i.north) -- ([xshift=2.2mm]t\i.south);}
  \node[anchor=east, font=\sffamily\small] at (15,7) {1\R{辺}{ぺん}（m）};
  \node[anchor=east, font=\sffamily\small] at (15,-7) {\R{面積}{めん|せき}（\unit{m^2}）};
  \node[anchor=west, font=\sffamily\footnotesize, text=bkTheme!80!black] at (104,2.4) {↓2\R{乗}{じょう}する};
  \node[anchor=west, font=\sffamily\footnotesize, text=bkBlue!80!black] at (104,-2.4) {↑もとにもどす};
\end{tikzpicture}
\end{bkfigure}

\noindent\begin{minipage}[c]{0.62\linewidth}
たとえば、\R{面積}{めん|せき}が\qty{9}{m^2}の\R{正方形}{せい|ほう|けい}なら、1\R{辺}{ぺん}\qty{1}{m}の\R{正方形}{せい|ほう|けい}が9こ\R{入}{はい}ります。9こを\R{正方形}{せい|ほう|けい}にならべるには、たてに3こ、よこに3こ（$3\times3=9$）です。だから1\R{辺}{ぺん}は\qty{3}{m}。\R{面積}{めん|せき}が\qty{25}{m^2}なら、$5\times5=25$ なので\qty{5}{m}です。

つまり、「2\R{乗}{じょう}すると9になる\R{数}{かず}は？」と\R{考}{かんが}えればよいのです。2\R{乗}{じょう}の\R{計算}{けい|さん}を、\R{逆}{ぎゃく}の\R{向}{む}きにたどっています。
\end{minipage}\hfill
\begin{minipage}[c]{0.35\linewidth}\centering
\begin{tikzpicture}[x=1mm, y=1mm]
  \foreach \i in {0,1,2} \foreach \j in {0,1,2} \filldraw[fill=bkGreen!22, draw=bkGreen!70!black, line width=0.5pt] (\i*8,\j*8) rectangle ++(8,8);
  \node[anchor=north, font=\sffamily\footnotesize, inner sep=0pt] at (12,-1.8) {\qty{3}{m}};
  \node[anchor=east, font=\sffamily\footnotesize, inner sep=0pt] at (-1.8,12) {\qty{3}{m}};
  \node[anchor=west, font=\sffamily\footnotesize, align=left, inner sep=0pt] at (26.5,12) {\qty{1}{m^2}が\\9こ};
\end{tikzpicture}
\end{minipage}

\bksubsection{2}{2\R{乗}{じょう}ではさんで、\R{少}{すこ}しずつ\R{近}{ちか}づく}

\leavevmode\label{bk:2-8:hasami}では、\bkhatenaref の\qty{2}{m^2}は？　\R{整数}{せい|すう}では、2\R{乗}{じょう}してちょうど2になる\R{数}{かず}はありません。$1^2=1$ は2より\R{小}{ちい}さく、$2^2=4$ は2より\R{大}{おお}きいので、1\R{辺}{ぺん}は1と2の\R{間}{あいだ}です。

\bkiiihead{\R{小数}{しょう|すう}の2\R{乗}{じょう}の\R{計算}{けい|さん}}
1.4を2\R{乗}{じょう}するには、どうすればよいでしょう。1.4は $\dfrac{14}{10}$ です（2-4）。\R{分数}{ぶん|すう}のかけ\R{算}{ざん}は、\R{分子}{ぶん|し}どうし・\R{分母}{ぶん|ぼ}どうしをかけるので、
\[ 1.4\times1.4=\frac{14}{10}\times\frac{14}{10}=\frac{14\times14}{10\times10}=\frac{196}{100}=1.96 \]
\R{小数点}{しょう|すう|てん}より\R{下}{した}のけたが1つずつの\R{数}{かず}をかけると、\R{答}{こた}えは\R{小数点}{しょう|すう|てん}より\R{下}{した}が2けたになります。\R{電卓}{でん|たく}で $1.4\times1.4$ をおしても、1.96と\R{出}{で}ます。

\bkiiihead{\R{表}{ひょう}で、2をはさむ}
1.1から0.1ずつ\R{大}{おお}きくして、2\R{乗}{じょう}を\R{表}{ひょう}にします。

\begin{bkcenter}
\begin{bktable}{colspec={Q[l,m,wd=20mm] *{5}{Q[c,m,wd=14mm]}}, rowsep=1.5pt, colsep=3pt, cells={font=\small}, column{1}={font=\sffamily\small}, cell{2}{5}={bg=bkBlue!18}, cell{2}{6}={bg=bkOrange!22}}
1\R{辺}{ぺん}（m） & 1.1 & 1.2 & 1.3 & 1.4 & 1.5 \\
2\R{乗}{じょう}（\unit{m^2}） & 1.21 & 1.44 & 1.69 & 1.96 & 2.25 \\
\end{bktable}
\bktabcaption{1.4と1.5のあいだで、2\R{乗}{じょう}が2をこえる}
\end{bkcenter}

$1.4^2=1.96$ は2より\R{小}{ちい}さく、$1.5^2=2.25$ は2より\R{大}{おお}きい。だから1\R{辺}{ぺん}は1.4と1.5の\R{間}{あいだ}です。こうして、2\R{乗}{じょう}が2より\R{小}{ちい}さい\R{数}{かず}と\R{大}{おお}きい\R{数}{かず}ではさみ、はばを$\dfrac{1}{10}$ずつせまくしていきます。\R{図}{ず}の1\R{本}{ぽん}めの\R{数直線}{すう|ちょく|せん}の1.4〜1.5の\R{間}{あいだ}を10に\R{分}{わ}けて\R{大}{おお}きくしたのが2\R{本}{ほん}め、その1.41〜1.42を\R{大}{おお}きくしたのが3\R{本}{ぼん}めです。

\begin{bkfigure}
\begin{tikzpicture}[x=1mm, y=1mm]
  \fill[bkTheme!9] (32,-1.4) -- (40,-1.4) -- (80,-11) -- (0,-11) -- cycle;
  \fill[bkTheme!9] (8,-19.4) -- (16,-19.4) -- (80,-29) -- (0,-29) -- cycle;
  \fill[bkTheme!45] (32,-1.4) rectangle (40,1.4);
  \bkiiinumline{(0,0)}{80}{1}{2}{10}
  \node[anchor=south east, font=\sffamily\footnotesize, inner sep=0pt] at (31.5,2.2) {1.4};
  \node[anchor=south west, font=\sffamily\footnotesize, inner sep=0pt] at (40.5,2.2) {1.5};
  \node[anchor=west, font=\footnotesize] at (84,0) {$1.4^2=1.96$、$1.5^2=2.25$};
  \begin{scope}[shift={(0,-18)}]
  \fill[bkTheme!45] (8,-1.4) rectangle (16,1.4);
  \bkiiinumline{(0,0)}{80}{1.4}{1.5}{10}
  \node[anchor=south east, font=\sffamily\footnotesize, inner sep=0pt] at (7.5,2.2) {1.41};
  \node[anchor=south west, font=\sffamily\footnotesize, inner sep=0pt] at (16.5,2.2) {1.42};
  \node[anchor=west, font=\footnotesize] at (84,0) {$1.41^2=1.9881$、$1.42^2=2.0164$};
  \end{scope}
  \begin{scope}[shift={(0,-36)}]
  \fill[bkTheme!45] (32,-1.4) rectangle (40,1.4);
  \bkiiinumline{(0,0)}{80}{1.41}{1.42}{10}
  \node[anchor=south east, font=\sffamily\footnotesize, inner sep=0pt] at (31.5,2.2) {1.414};
  \node[anchor=south west, font=\sffamily\footnotesize, inner sep=0pt] at (40.5,2.2) {1.415};
  \node[anchor=west, font=\footnotesize, align=left] at (84,0) {$1.414^2=1.999396$\\$1.415^2=2.002225$};
  \end{scope}
\end{tikzpicture}
\bkfigcaption{2\R{乗}{じょう}すると2になる\R{数}{かず}を、はさみうちでさがす。\R{色}{いろ}の\R{部分}{ぶ|ぶん}を\R{次}{つぎ}の\R{線}{せん}で10\R{倍}{ばい}に\R{大}{おお}きくしている。}
\end{bkfigure}

\bkhatenaref の\R{答}{こた}えは\R{約}{やく}\qty{1.41}{m}です。\qty{1}{m}は\qty{100}{cm}なので、0.41\,mは $0.41\times100=41$ で\qty{41}{cm}。\R{約}{やく}1\,m\,41\,cmです。でも、$1.414^2$ も、まだぴったり2ではありません。どこまでせまくしても、ぴったりにはならないのです（2-9）。

\bksubsection{3}{\R{記号}{き|ごう}$\sqrt{\phantom{2}}$（ルート）}

\leavevmode\label{bk:2-8:root}2\R{乗}{じょう}すると2になる\R{正}{せい}の\R{数}{かず}を、$\sqrt{2}$ と\R{書}{か}き、「ルート2」と\R{読}{よ}みます。\bkidx{るーと}{$\sqrt{\phantom{2}}$（ルート）}
\[ \sqrt{2}=1.41421356\cdots,\qquad \sqrt{9}=3,\qquad \sqrt{25}=5 \]

\noindent\begin{minipage}[c]{0.64\linewidth}
2\R{乗}{じょう}して9になる\R{数}{かず}は、3だけではありません。$(-3)^2=(-3)\times(-3)=9$ だからです（2-2：マイナス×マイナスはプラス）。2\R{乗}{じょう}して□になる\R{数}{かず}を、□の\textbf{\R{平方根}{へい|ほう|こん}}といいます。\bkidx{へいほうこん}{\R{平方根}{へい|ほう|こん}}9の\R{平方根}{へい|ほう|こん}は3と$-3$の2つ。$\sqrt{9}$ は、そのうちの\R{正}{せい}のほうの3です。\R{長}{なが}さはマイナスにならないので、1\R{辺}{ぺん}を\R{求}{もと}めるときは\R{正}{せい}のほうを\R{使}{つか}います。
\end{minipage}\hfill
\begin{minipage}[c]{0.33\linewidth}\centering
\begin{tikzpicture}[x=1mm, y=1mm]
  \bkdisplay{(0,0)}{38}{1.414213562}
  \foreach \k/\t in {0/2,1/{\sqrt{\ }}} \node[draw=bkInk!60, fill=white, rounded corners=1mm, minimum width=9mm, minimum height=6.5mm, font=\small] at (12+\k*14,-6.5) {$\t$};
  \node[anchor=north, font=\sffamily\footnotesize] at (19,-11) {2を\R{入}{い}れて $\sqrt{\ }$ キー};
\end{tikzpicture}
\end{minipage}

\noindent\textbf{0とマイナスの\R{数}{かず}は？}\quad $0^2=0$ なので、0の\R{平方根}{へい|ほう|こん}は0だけです。また、\R{正}{せい}の\R{数}{かず}を2\R{乗}{じょう}しても、\R{負}{ふ}の\R{数}{かず}を2\R{乗}{じょう}しても、\R{答}{こた}えはプラスになります。だから、この\R{巻}{かん}で\R{学}{まな}ぶ\R{数}{かず}の\R{中}{なか}には、2\R{乗}{じょう}して $-4$ のようなマイナスになる\R{数}{かず}はありません（そういう\R{数}{かず}は6-8で\R{新}{あたら}しく\R{考}{かんが}えます）。

\noindent\textbf{$\sqrt{2}\times\sqrt{2}=2$}\quad $\sqrt{2}$ は「2\R{乗}{じょう}すると2になる\R{数}{かず}」なので、2\R{回}{かい}かけると、\R{決}{き}めたとおりに2になります。\R{電卓}{でん|たく}で $1.41421356\times1.41421356$ を\R{計算}{けい|さん}すると、1.9999999…と、2にとても\R{近}{ちか}い\R{数}{かず}が\R{出}{で}ます。

\bksubsection{4}{$\sqrt{\phantom{2}}$ どうしのかけ\R{算}{ざん}}

\leavevmode\label{bk:2-8:kake}$\sqrt{2}\times\sqrt{3}$ を\R{電卓}{でん|たく}で\R{計算}{けい|さん}すると、$1.41421\cdots\times1.73205\cdots=2.44948\cdots$。これは $\sqrt{6}$ と\R{同}{おな}じ\R{数}{かず}です。なぜでしょう。$\sqrt{2}\times\sqrt{3}$ を2\R{乗}{じょう}してみます。
\[ (\sqrt{2}\times\sqrt{3})\times(\sqrt{2}\times\sqrt{3})=(\sqrt{2}\times\sqrt{2})\times(\sqrt{3}\times\sqrt{3})=2\times3=6 \]
1\R{行}{ぎょう}めから2\R{行}{ぎょう}めへは、かけ\R{算}{ざん}の\R{順番}{じゅん|ばん}を\R{入}{い}れかえただけです（かけ\R{算}{ざん}は、どの\R{順}{じゅん}にかけても\R{答}{こた}えが\R{同}{おな}じ）。$\sqrt{2}\times\sqrt{3}$ は、2\R{乗}{じょう}して6になる\R{正}{せい}の\R{数}{かず}なので、$\sqrt{2}\times\sqrt{3}=\sqrt{6}$ です。$\sqrt{\phantom{2}}$ どうしのかけ\R{算}{ざん}は、\R{中}{なか}の\R{数}{かず}どうしをかけてよいのです。たとえば $\sqrt{2}\times\sqrt{8}=\sqrt{16}=4$。

\bksubsection{5}{\R{振}{ふ}り\R{子}{こ}の\R{周期}{しゅう|き}は、\R{長}{なが}さの $\sqrt{\phantom{2}}$ に\R{比例}{ひ|れい}する}

\leavevmode\label{bk:2-8:furiko}\bkyattemiyouref の\R{表}{ひょう}では、\R{糸}{いと}の\R{長}{なが}さを4\R{倍}{ばい}にすると\R{周期}{しゅう|き}は2\R{倍}{ばい}、9\R{倍}{ばい}にすると3\R{倍}{ばい}になりました。$4=2^2$、$9=3^2$ です。つまり、\R{周期}{しゅう|き}の\R{倍}{ばい}は、\R{長}{なが}さの\R{倍}{ばい}の $\sqrt{\phantom{2}}$ になっています（$\sqrt{4}=2$、$\sqrt{9}=3$）。

\noindent\begin{minipage}[c]{0.5\linewidth}
では、\R{長}{なが}さを2\R{倍}{ばい}の\qty{20}{cm}にしたら？　\R{予想}{よ|そう}は $\sqrt{2}$\R{倍}{ばい}で、$0.63\times1.41=0.89$\R{秒}{びょう}。はかってみると\R{約}{やく}0.90\R{秒}{びょう}で、\R{予想}{よ|そう}どおりです。グラフは\R{直線}{ちょく|せん}ではなく、だんだんねていく\R{曲線}{きょく|せん}になります。
\end{minipage}\hfill
\begin{minipage}[c]{0.48\linewidth}\centering
\begin{tikzpicture}
\begin{axis}[bkplot, width=52mm, height=42mm, xmin=0, xmax=1.08, ymin=0, ymax=2.35, xtick={0,0.2,...,1}, ytick={0,0.5,...,2},
  xticklabels={0,0.2,0.4,0.6,0.8,1}, yticklabels={0,0.5,1,1.5,2},
  xlabel={\R{長}{なが}さ（m）}, ylabel={\R{周期}{しゅう|き}（\R{秒}{びょう}）}]
  \addplot[bkcurve, domain=0:1.0] {2.007*sqrt(x)};
  \addplot[only marks, mark=*, mark size=1.6pt, bkOrange] coordinates {(0.1,0.63) (0.2,0.90) (0.4,1.27) (0.9,1.90)};
\end{axis}
\end{tikzpicture}
\end{minipage}

\R{一方}{いっ|ぽう}が2\R{倍}{ばい}、3\R{倍}{ばい}、……になると、もう\R{一方}{いっ|ぽう}も2\R{倍}{ばい}、3\R{倍}{ばい}、……になるとき、\textbf{\R{比例}{ひ|れい}する}といいます（くわしくは2-15）。\bkidx{ふりこのしゅうき}{\R{振}{ふ}り\R{子}{こ}の\R{周期}{しゅう|き}}\R{振}{ふ}り\R{子}{こ}の\R{周期}{しゅう|き}は、\R{長}{なが}さに\R{比例}{ひ|れい}するのではなく、\textbf{\R{糸}{いと}の\R{長}{なが}さの $\sqrt{\phantom{2}}$ に\R{比例}{ひ|れい}します}。

\bkiiihead{\R{表}{ひょう}で\R{確}{たし}かめる：\R{周期}{しゅう|き}÷$\sqrt{\text{\R{長}{なが}さ}}$ はいつも\R{約}{やく}2}
\R{比例}{ひ|れい}しているなら、\R{周期}{しゅう|き}を $\sqrt{\phantom{2}}$（\R{長}{なが}さ）でわった\R{数}{かず}は、いつも\R{同}{おな}じになるはずです。\R{長}{なが}さをmにして、\R{電卓}{でん|たく}で\R{計算}{けい|さん}してみます（\qty{10}{cm}は $10\div100=0.1$ で\qty{0.1}{m}）。

\begin{bkcenter}
\begin{bktable}{colspec={Q[l,m,wd=36mm] *{4}{Q[c,m,wd=17mm]}}, rowsep=1.5pt, colsep=3pt, cells={font=\small}, column{1}={font=\sffamily\small}, row{4}={bg=bkTheme!10}}
\R{長}{なが}さ（m） & 0.1 & 0.2 & 0.4 & 0.9 \\
$\sqrt{\text{\R{長}{なが}さ}}$ & 0.316 & 0.447 & 0.632 & 0.949 \\
\R{周期}{しゅう|き}（\R{秒}{びょう}） & 0.63 & 0.90 & 1.27 & 1.90 \\
\R{周期}{しゅう|き}÷$\sqrt{\text{\R{長}{なが}さ}}$ & 1.99 & 2.01 & 2.01 & 2.00 \\
\end{bktable}
\bktabcaption{どの\R{長}{なが}さでも、ほぼ2になる}
\end{bkcenter}

どれもほぼ2です。だから、\textbf{\R{長}{なが}さ（m）の $\sqrt{\phantom{2}}$ を2\R{倍}{ばい}すると、\R{周期}{しゅう|き}（\R{秒}{びょう}）のおよその\R{値}{あたい}}になります。この2は、はかった\R{表}{ひょう}から\R{見}{み}つけた\R{数}{かず}です。\qty{1}{m}の\R{振}{ふ}り\R{子}{こ}なら $\sqrt{1}=1$ で、\R{周期}{しゅう|き}は\R{約}{やく}2\R{秒}{びょう}です。なぜ $\sqrt{\phantom{2}}$ になるのか、なぜ\R{約}{やく}2なのかは、6-6で\R{式}{しき}から\R{求}{もと}めます。

\noindent\textbf{グラフが\R{曲}{ま}がるわけ}\quad \R{長}{なが}さを0.1から0.4へ0.3ふやすと、\R{周期}{しゅう|き}は0.64\R{秒}{びょう}ふえます。でも0.4から0.9へ、もっと\R{多}{おお}い0.5ふやしても、\R{周期}{しゅう|き}は0.63\R{秒}{びょう}しかふえません。\R{長}{なが}くなるほど、\R{周期}{しゅう|き}はふえにくくなるので、グラフはだんだんねていきます。

\begin{motto}[2]
\par\addvspace{3pt}\textbf{\R{電卓}{でん|たく}の $\sqrt{\phantom{2}}$ キーがないとき}\par\noindent 2を1.5でわると1.333…。1.5と1.333…の\R{平均}{へい|きん}は1.4167。こんどは2を1.4167でわって、また\R{平均}{へい|きん}をとると1.414216。2\R{回}{かい}くり\R{返}{かえ}すだけで、$\sqrt{2}=1.414213\cdots$ のすぐ\R{近}{ちか}くまで\R{来}{き}ました。コンピュータも、にた\R{方法}{ほう|ほう}で $\sqrt{\phantom{2}}$ を\R{計算}{けい|さん}しています。
\end{motto}

\begin{taisetsu}
2\R{乗}{じょう}すると□になる\R{数}{かず}を、□の\R{平方根}{へい|ほう|こん}という。\\（□が0より\R{大}{おお}きいとき）\R{正}{せい}と\R{負}{ふ}の2つある。
\R{正}{せい}のほうを $\sqrt{\text{□}}$ と\R{書}{か}く。\quad $\sqrt{9}=3$、\ $\sqrt{2}=1.41421\cdots$\\
$\sqrt{2}\times\sqrt{3}=\sqrt{6}$\quad（かけ\R{算}{ざん}は\R{中}{なか}どうしをかけてよい）\\
\R{振}{ふ}り\R{子}{こ}の\R{周期}{しゅう|き}は、\R{糸}{いと}の\R{長}{なが}さの $\sqrt{\phantom{2}}$ に\R{比例}{ひ|れい}する。
\end{taisetsu}

\begin{kotoba}
\small
\bkword{\R{平方根}{へい|ほう|こん}}{2\R{乗}{じょう}すると、その\R{数}{かず}になる\R{数}{かず}。9の\R{平方根}{へい|ほう|こん}は3と$-3$。0の\R{平方根}{へい|ほう|こん}は0だけ。}
\bkword{\R{周期}{しゅう|き}}{\R{振}{ふ}り\R{子}{こ}が1\R{往復}{おう|ふく}する\R{時間}{じ|かん}。くり\R{返}{かえ}す\R{動}{うご}きの1\R{回分}{かい|ぶん}の\R{時間}{じ|かん}。}
\bkword{\R{比例}{ひ|れい}}{\R{一方}{いっ|ぽう}が2\R{倍}{ばい}、3\R{倍}{ばい}になると、もう\R{一方}{いっ|ぽう}も2\R{倍}{ばい}、3\R{倍}{ばい}になること（2-15）。}
\end{kotoba}

\bkkigocards{\bkkigocard{$\sqrt{\phantom{00}}$}{ルート}{2\R{乗}{じょう}すると\R{中}{なか}の\R{数}{かず}になる\R{正}{せい}の\R{数}{かず}。$\sqrt{9}=3$}{square root}{2-8}}{%
\begin{tikzpicture}[x=1mm, y=1mm]
  \filldraw[fill=bkTheme!14, draw=bkTheme!80!black, line width=0.7pt] (0,0) -- (16,16) -- (0,32) -- (-16,16) -- cycle;
  \filldraw[fill=bkBlue!25, fill opacity=0.6, draw=bkBlue!70!black, line width=0.7pt] (0,0) rectangle (16,16);
  \draw[bkTheme!80!black, line width=0.5pt] (0,16) -- (0,32) (-16,16) -- (0,16);
  \draw[bkTheme!85!black, line width=1.6pt] (0,0) -- (16,16);
  \node[font=\small, text=bkTheme!85!black, inner sep=0pt] at (11.6,4.3) {$\sqrt{2}$};
  \node[anchor=north, font=\sffamily\footnotesize, inner sep=0pt] at (8,-1.5) {1};
  \node[anchor=west, font=\sffamily\footnotesize, text=bkBlue!70!black] at (17.5,8) {\R{面積}{めん|せき}1};
  \node[anchor=south east, font=\sffamily\footnotesize, text=bkTheme!80!black, inner sep=0pt] at (-10,25) {\R{面積}{めん|せき}2};
  \node[anchor=north, font=\sffamily\footnotesize, align=center, inner sep=0pt] at (2,-6.5) {\R{正方形}{せい|ほう|けい}の\R{対角線}{たい|かく|せん}が $\sqrt{2}$};
\end{tikzpicture}}

%% ---------------- 例題 ----------------
\bksection{reidai}{\R{例題}{れい|だい}でたしかめよう}

\begin{reidai}[\R{平方根}{へい|ほう|こん}を\R{求}{もと}める]
（1）49の\R{平方根}{へい|ほう|こん}を\R{求}{もと}めましょう。（2）$\sqrt{81}$\quad（3）$\sqrt{0.09}$ を\R{求}{もと}めましょう。\par
\kangaekata 2\R{乗}{じょう}すると、その\R{数}{かず}になる\R{数}{かず}をさがします。\R{平方根}{へい|ほう|こん}は\R{正}{せい}と\R{負}{ふ}の2つ、$\sqrt{\phantom{2}}$ は\R{正}{せい}のほうです。\R{小数}{しょう|すう}のときは、0.1、0.2、0.3、……を\R{順}{じゅん}に2\R{乗}{じょう}してさがします。\par
\shiki （1）$7^2=49$、\ $(-7)^2=49$\par
\hspace*{2.6\zw}（2）$9^2=81$\par
\hspace*{2.6\zw}（3）$0.3^2=\dfrac{3}{10}\times\dfrac{3}{10}=\dfrac{9}{100}=0.09$\par
\kotae （1）7と$-7$\quad（2）9\quad（3）0.3
\end{reidai}

\begin{jibun}
（1）64の\R{平方根}{へい|ほう|こん}を\R{求}{もと}めましょう。（2）$\sqrt{100}$\quad（3）$\sqrt{0.16}$ を\R{求}{もと}めましょう。
\bkjans{（1）\nobreak 8と$-8$\quad（2）\nobreak 10\quad（3）\nobreak 0.4}
\end{jibun}

\begin{reidai}[はさみうちで1\R{辺}{ぺん}を\R{求}{もと}める]
\R{面積}{めん|せき}が\qty{10}{m^2}の\R{正方形}{せい|ほう|けい}の1\R{辺}{ぺん}を、\R{小数}{しょう|すう}\R{第}{だい}2\R{位}{い}まではさんで\R{求}{もと}め、\R{四捨五入}{し|しゃ|ご|にゅう}して\R{小数}{しょう|すう}\R{第}{だい}1\R{位}{い}までのがい\R{数}{すう}で\R{答}{こた}えましょう。\par
\kangaekata 2\R{乗}{じょう}が10より\R{小}{ちい}さい\R{数}{かず}と\R{大}{おお}きい\R{数}{かず}ではさみます。\par
\shiki $3^2=9$、$4^2=16$ なので、3と4の\R{間}{あいだ}\par
\hspace*{2.6\zw}$3.1^2=9.61$、$3.2^2=10.24$ なので、3.1と3.2の\R{間}{あいだ}\par
\hspace*{2.6\zw}$3.16^2=9.9856$、$3.17^2=10.0489$ なので、3.16と3.17の\R{間}{あいだ}\par
\kotae 3.16…を\R{四捨五入}{し|しゃ|ご|にゅう}して、\R{約}{やく}\qty{3.2}{m}
\end{reidai}

\begin{jibun}
\R{面積}{めん|せき}が\qty{5}{cm^2}の\R{正方形}{せい|ほう|けい}の1\R{辺}{ぺん}を、\R{小数}{しょう|すう}\R{第}{だい}2\R{位}{い}まではさんで\R{求}{もと}め、\R{小数}{しょう|すう}\R{第}{だい}1\R{位}{い}までのがい\R{数}{すう}で\R{答}{こた}えましょう。
\bkjans{$2.2^2=4.84$、$2.3^2=5.29$、$2.23^2=4.9729$、$2.24^2=5.0176$ より 2.23…で、\R{約}{やく}\qty{2.2}{cm}}
\end{jibun}

\begin{reidai}[$\sqrt{\phantom{2}}$ の\R{大}{おお}きさの\R{見当}{けん|とう}をつける]
$\sqrt{30}$ は、どの2つの\R{整数}{せい|すう}の\R{間}{あいだ}にありますか。\par
\kangaekata \R{整数}{せい|すう}を\R{順}{じゅん}に2\R{乗}{じょう}して、30より\R{小}{ちい}さいものと\R{大}{おお}きいものではさみます。\par
\shiki $5^2=25$ は30より\R{小}{ちい}さく、$6^2=36$ は30より\R{大}{おお}きい\par
\hspace*{2.6\zw}2\R{乗}{じょう}して30になる\R{数}{かず}は、5より\R{大}{おお}きく6より\R{小}{ちい}さい\par
\kotae 5と6の\R{間}{あいだ}
\end{reidai}

\begin{jibun}
$\sqrt{60}$ は、どの2つの\R{整数}{せい|すう}の\R{間}{あいだ}にありますか。
\bkjans{$7^2=49$、$8^2=64$ より、7と8の\R{間}{あいだ}}
\end{jibun}

\begin{reidai}[\R{電卓}{でん|たく}と $\sqrt{\phantom{2}}$ のかけ\R{算}{ざん}]
（1）$\sqrt{3}$ を\R{電卓}{でん|たく}で\R{求}{もと}め、\R{小数}{しょう|すう}\R{第}{だい}2\R{位}{い}までのがい\R{数}{すう}にしましょう。（2）$\sqrt{2}\times\sqrt{8}$\quad（3）$\sqrt{3}\times\sqrt{12}$ を\R{計算}{けい|さん}しましょう。\par
\kangaekata かけ\R{算}{ざん}は、\R{中}{なか}の\R{数}{かず}どうしをかけます。\par
\shiki （1）$\sqrt{3}=1.7320\cdots$ の\R{小数}{しょう|すう}\R{第}{だい}3\R{位}{い}を\R{四捨五入}{し|しゃ|ご|にゅう}\par
\hspace*{2.6\zw}（2）$\sqrt{2}\times\sqrt{8}=\sqrt{16}=4$\qquad（3）$\sqrt{3}\times\sqrt{12}=\sqrt{36}=6$\par
\kotae （1）1.73\quad（2）4\quad（3）6
\end{reidai}

\begin{jibun}
（1）$\sqrt{5}$ を\R{電卓}{でん|たく}で\R{求}{もと}め、\R{小数}{しょう|すう}\R{第}{だい}2\R{位}{い}までのがい\R{数}{すう}にしましょう。（2）$\sqrt{2}\times\sqrt{18}$\quad（3）$\sqrt{5}\times\sqrt{5}$ を\R{計算}{けい|さん}しましょう。
\bkjans{（1）\nobreak 2.24（$2.2360\cdots$）\quad（2）\nobreak $\sqrt{36}=6$\quad（3）\nobreak $\sqrt{25}=5$}
\end{jibun}

\begin{column}{1メートルは\R{振}{ふ}り\R{子}{こ}で\R{決}{き}まるはずだった？}
230\R{年}{ねん}ほど\R{前}{まえ}、フランスで\R{長}{なが}さの\R{単位}{たん|い}を\R{決}{き}めるとき、「\R{片道}{かた|みち}がちょうど1\R{秒}{びょう}の\R{振}{ふ}り\R{子}{こ}の\R{長}{なが}さ」（\R{約}{やく}\qty{0.994}{m}）を1メートルにしよう、という\R{案}{あん}がありました。でも、\R{周期}{しゅう|き}は\R{場所}{ば|しょ}によって\R{少}{すこ}し\R{変}{か}わるので、\R{地球}{ち|きゅう}の\R{大}{おお}きさをもとに\R{決}{き}めることになりました。それでも、2つの\R{長}{なが}さはほとんど\R{同}{おな}じです（2-10のコラムも\R{見}{み}てみよう）。
\end{column}

\begin{reidai}[\R{面積}{めん|せき}を\R{何倍}{なん|ばい}かにすると、1\R{辺}{ぺん}は？]
1\R{辺}{ぺん}\qty{1}{m}の\R{正方形}{せい|ほう|けい}があります。（1）\R{面積}{めん|せき}を4\R{倍}{ばい}にするには、1\R{辺}{ぺん}を\R{何倍}{なん|ばい}にすればよいですか。（2）\R{面積}{めん|せき}を2\R{倍}{ばい}にするには、\R{約}{やく}\R{何倍}{なん|ばい}ですか。\par
\kangaekata 1\R{辺}{ぺん}を□\R{倍}{ばい}にすると、たても□\R{倍}{ばい}、よこも□\R{倍}{ばい}になるので、\R{面積}{めん|せき}は□×□\R{倍}{ばい}になります（\R{下}{した}の\R{図}{ず}）。□×□が4や2になる□をさがします。\par
\begin{center}
\begin{tikzpicture}[x=1mm, y=1mm]
  \filldraw[fill=bkBlue!25, draw=bkBlue!70!black, line width=0.6pt] (0,0) rectangle (9,9);
  \node[anchor=north, font=\sffamily\footnotesize, inner sep=0pt] at (4.5,-1.6) {1\R{辺}{ぺん}1};
  \foreach \i in {0,1} \foreach \j in {0,1} \filldraw[fill=bkOrange!25, draw=bkOrange!80!black, line width=0.6pt] (16+\i*9,\j*9) rectangle ++(9,9);
  \node[anchor=north, font=\sffamily\footnotesize, inner sep=0pt] at (25,-1.6) {1\R{辺}{ぺん}2};
  \node[anchor=south, font=\sffamily\footnotesize, inner sep=0pt] at (25,19.6) {\R{面積}{めん|せき}は4\R{倍}{ばい}};
  \node[anchor=west, font=\sffamily\footnotesize, align=left, inner sep=0pt] at (40,9) {1\R{辺}{ぺん}を2\R{倍}{ばい}にすると、\\たても2\R{倍}{ばい}、よこも2\R{倍}{ばい}。\\\R{面積}{めん|せき}は $2\times2=4$（\R{倍}{ばい}）};
\end{tikzpicture}
\end{center}
\shiki （1）$2\times2=4$ なので、1\R{辺}{ぺん}は $\sqrt{4}=2$ で2\R{倍}{ばい}\par
\hspace*{2.6\zw}（2）□×□＝2 になる□は $\sqrt{2}=1.41\cdots$\par
\kotae （1）2\R{倍}{ばい}\quad（2）\R{約}{やく}1.41\R{倍}{ばい}
\end{reidai}

\begin{jibun}
\R{正方形}{せい|ほう|けい}の（1）\R{面積}{めん|せき}を16\R{倍}{ばい}にするには、1\R{辺}{ぺん}を\R{何倍}{なん|ばい}に？（2）\R{面積}{めん|せき}を3\R{倍}{ばい}にするには、1\R{辺}{ぺん}を\R{約}{やく}\R{何倍}{なん|ばい}に？（\R{電卓}{でん|たく}を\R{使}{つか}い、\R{小数第}{しょう|すう|だい}2\R{位}{い}までのがい\R{数}{すう}で）
\bkjans{（1）\nobreak $\sqrt{16}=4$ で4\R{倍}{ばい}\quad（2）\nobreak $\sqrt{3}=1.732\cdots$ で\R{約}{やく}1.73\R{倍}{ばい}}
\end{jibun}

\begin{reidai}[\R{振}{ふ}り\R{子}{こ}の\R{周期}{しゅう|き}]
\R{長}{なが}さ\qty{25}{cm}の\R{振}{ふ}り\R{子}{こ}の\R{周期}{しゅう|き}は\R{約}{やく}1.0\R{秒}{びょう}です。（1）\R{長}{なが}さを\qty{1}{m}にすると、\R{周期}{しゅう|き}は\R{何秒}{なん|びょう}ですか。（2）\R{周期}{しゅう|き}を3.0\R{秒}{びょう}にするには、\R{長}{なが}さを\R{何}{なん}mにすればよいですか。\par
\kangaekata \R{周期}{しゅう|き}の\R{倍}{ばい}は、\R{長}{なが}さの\R{倍}{ばい}の $\sqrt{\phantom{2}}$。\R{逆}{ぎゃく}に、\R{長}{なが}さの\R{倍}{ばい}は、\R{周期}{しゅう|き}の\R{倍}{ばい}の2\R{乗}{じょう}です（$\sqrt{\text{□}}$ が3なら、□は $3^2=9$）。\R{倍}{ばい}をくらべるときは、\R{単位}{たん|い}をそろえます。\par
\shiki （1）① \qty{1}{m}は\qty{100}{cm}なので、\R{長}{なが}さは $100\div25=4$（\R{倍}{ばい}）\par
\hspace*{2.6\zw}\phantom{（1）}② \R{周期}{しゅう|き}は $\sqrt{4}=2$（\R{倍}{ばい}）なので、$1.0\times2=2.0$\par
\hspace*{2.6\zw}（2）① \R{周期}{しゅう|き}は $3.0\div1.0=3$（\R{倍}{ばい}）\par
\hspace*{2.6\zw}\phantom{（2）}② \R{長}{なが}さは $3^2=9$（\R{倍}{ばい}）なので、$25\times9=225$（cm）\par
\hspace*{2.6\zw}\phantom{（2）}③ \qty{100}{cm}が\qty{1}{m}なので、$225\div100=2.25$（m）\par
\kotae （1）\R{約}{やく}2.0\R{秒}{びょう}\quad（2）\qty{2.25}{m}
\end{reidai}

\begin{jibun}
\R{長}{なが}さ\qty{10}{cm}の\R{振}{ふ}り\R{子}{こ}の\R{周期}{しゅう|き}は0.63\R{秒}{びょう}です。（1）\R{長}{なが}さ\qty{1.6}{m}のときの\R{周期}{しゅう|き}は？（2）\R{周期}{しゅう|き}を1.89\R{秒}{びょう}にするには、\R{長}{なが}さを\R{何}{なん}cmに？
\bkjans{（1）\nobreak \qty{1.6}{m}＝\qty{160}{cm}で16\R{倍}{ばい}。$\sqrt{16}=4$\R{倍}{ばい}で、$0.63\times4=2.52$ で2.52\R{秒}{びょう}\quad（2）\nobreak $1.89\div0.63=3$\R{倍}{ばい}。\R{長}{なが}さは $3^2=9$\R{倍}{ばい}で\qty{90}{cm}}
\end{jibun}

%% ---------------- まちがえやすいところ・物理工学 ----------------
\bksection{tsumazuki}{まちがえやすいところをチェック}

\begin{tsumazuki}{$\sqrt{\phantom{2}}$ のたし\R{算}{ざん}で、\R{中}{なか}どうしをたしてしまう}
\bkng{$\sqrt{2}+\sqrt{3}=\sqrt{5}$}
\bkok{$\sqrt{2}+\sqrt{3}=1.414\cdots+1.732\cdots=3.146\cdots$}
$\sqrt{5}=2.236\cdots$ なので、ちがいます。\R{中}{なか}どうしをかけてよいのは、かけ\R{算}{ざん}のときだけです。
\end{tsumazuki}

\begin{tsumazuki}{2\R{乗}{じょう}を「2をかける」と\R{思}{おも}う}
\bkng{$1.4^2=1.4\times2=2.8$}
\bkok{$1.4^2=1.4\times1.4=1.96$}
2\R{乗}{じょう}は「\R{同}{おな}じ\R{数}{かず}を2\R{回}{かい}かける」ことです（2-7）。2をかけるのではありません。
\end{tsumazuki}

\begin{tsumazuki}{\R{平方根}{へい|ほう|こん}を1つだけ\R{答}{こた}える}
\bkng{9の\R{平方根}{へい|ほう|こん}は3}
\bkok{9の\R{平方根}{へい|ほう|こん}は3と$-3$。$\sqrt{9}$ は\R{正}{せい}のほうで3}
$(-3)^2=9$ なので、$-3$ も9の\R{平方根}{へい|ほう|こん}です。
\end{tsumazuki}

\begin{tsumazuki}{\R{長}{なが}さを2\R{倍}{ばい}にすると、\R{周期}{しゅう|き}も2\R{倍}{ばい}と\R{考}{かんが}える}
\bkng{\R{長}{なが}さ2\R{倍}{ばい}→\R{周期}{しゅう|き}2\R{倍}{ばい}}
\bkok{\R{長}{なが}さ2\R{倍}{ばい}→\R{周期}{しゅう|き} $\sqrt{2}$ \R{倍}{ばい}（\R{約}{やく}1.41\R{倍}{ばい}）}
\R{周期}{しゅう|き}は、\R{長}{なが}さの $\sqrt{\phantom{2}}$ に\R{比例}{ひ|れい}します。\R{周期}{しゅう|き}を2\R{倍}{ばい}にするには、\R{長}{なが}さを4\R{倍}{ばい}にします。
\end{tsumazuki}

\begin{butsuri}
\noindent\begin{minipage}[c]{0.62\linewidth}
\par\addvspace{3pt}\textbf{A\R{判}{ばん}の\R{紙}{かみ}}\par\noindent A4の\R{紙}{かみ}は\qty{210}{mm}×\qty{297}{mm}。$297\div210=1.414\cdots$ で、\R{短}{みじか}い\R{辺}{へん}と\R{長}{なが}い\R{辺}{へん}の\R{比}{ひ}は $1:\sqrt{2}$ です。この\R{形}{かたち}だと、\R{半分}{はん|ぶん}に\R{切}{き}っても\R{同}{おな}じ\R{形}{かたち}になります（A5は\qty{148}{mm}×\qty{210}{mm}、$210\div148=1.42$）。\R{面積}{めん|せき}が2\R{倍}{ばい}のA3にすると、\R{辺}{へん}は $\sqrt{2}$\R{倍}{ばい}。コピー\R{機}{き}でA4をA3に\R{広}{ひろ}げるときの「141\,\%」は、$\sqrt{2}=1.41$ のことです。
\end{minipage}\hfill
\begin{minipage}[c]{0.35\linewidth}\centering
\begin{tikzpicture}[x=1mm, y=1mm]
  \filldraw[fill=bkPurple!12, draw=bkPurple!70!black, line width=0.7pt] (0,0) rectangle (21.0,29.7);
  \filldraw[fill=bkPurple!30, draw=bkPurple!70!black, line width=0.7pt] (21.0,14.85) rectangle (42.0,29.7);
  \filldraw[fill=bkPurple!30, draw=bkPurple!70!black, line width=0.7pt] (21.0,0) rectangle (42.0,14.85);
  \node[font=\sffamily\small] at (10.5,14.85) {A4};
  \node[font=\sffamily\small] at (31.5,22.3) {A5};
  \node[font=\sffamily\small] at (31.5,7.4) {A5};
  \node[anchor=north, font=\sffamily\footnotesize] at (21,-1.5) {\R{全体}{ぜん|たい}がA3};
\end{tikzpicture}
\end{minipage}

\vspace{8pt}
\noindent\begin{minipage}[c]{0.75\linewidth}
\par\addvspace{3pt}\textbf{\R{振}{ふ}り\R{子}{こ}\R{時計}{どけい}}\par\noindent \R{長}{なが}さ\R{約}{やく}\qty{1}{m}の\R{振}{ふ}り\R{子}{こ}は、\R{周期}{しゅう|き}が\R{約}{やく}2\R{秒}{びょう}。\R{片道}{かた|みち}がちょうど1\R{秒}{びょう}です。\R{時計}{とけい}が\R{進}{すす}むときは、おもりを\R{下}{さ}げて\R{長}{なが}さを\R{長}{なが}くし、\R{周期}{しゅう|き}を\R{長}{なが}くします。\R{長}{なが}さの\R{変}{か}え\R{方}{かた}は、\R{周期}{しゅう|き}の\R{変}{か}え\R{方}{かた}の2\R{乗}{じょう}になることに\R{注意}{ちゅう|い}します。
\end{minipage}\hfill
\begin{minipage}[c]{0.22\linewidth}\centering
\begin{tikzpicture}[x=1mm, y=1mm]
  \filldraw[fill=bkStem!60, draw=bkStem!80!black, line width=0.6pt, rounded corners=1mm] (-8,-34) rectangle (8,6);
  \filldraw[fill=white, draw=bkInk!70, line width=0.6pt] (0,0) circle (5.2);
  \draw[bkInk, line width=0.8pt] (0,0) -- (0,3.6) (0,0) -- (2.6,-1);
  \draw[bkInk, line width=0.7pt] (0,-5.2) -- (3,-28);
  \fill[bkYellow!80!black] (3,-28) circle (3);
\end{tikzpicture}
\end{minipage}
\end{butsuri}



\clearpage
%% ---------------- 章末問題 ----------------
\bkendproblems
\bkproblemgroup{\R{基本}{き|ほん}}{2}

\bkproblem \R{次}{つぎ}の\R{数}{かず}を\R{求}{もと}めましょう。\quad（1）36の\R{平方根}{へい|ほう|こん}\quad（2）$\sqrt{121}$\quad（3）$\sqrt{0.49}$
\bkpans{（1）\nobreak 6と$-6$\quad（2）\nobreak 11\quad（3）\nobreak 0.7}

\bkproblem \R{面積}{めん|せき}が\qty{30}{m^2}の\R{正方形}{せい|ほう|けい}の1\R{辺}{ぺん}を、\R{小数}{しょう|すう}\R{第}{だい}2\R{位}{い}まではさんで\R{求}{もと}め、\R{小数}{しょう|すう}\R{第}{だい}1\R{位}{い}までのがい\R{数}{すう}で\R{答}{こた}えましょう。
\bkpans{$5.4^2=29.16$、$5.5^2=30.25$、$5.47^2=29.9209$、$5.48^2=30.0304$ より 5.47…で、\R{約}{やく}\qty{5.5}{m}}

\bkproblem \R{電卓}{でん|たく}で\R{求}{もと}め、\R{小数}{しょう|すう}\R{第}{だい}2\R{位}{い}までのがい\R{数}{すう}にしましょう。\quad（1）$\sqrt{7}$\quad（2）$\sqrt{11}$
\bkpans{（1）\nobreak 2.65\quad（2）\nobreak 3.32}

\bkproblem \R{次}{つぎ}の\R{計算}{けい|さん}をしましょう。\quad（1）$\sqrt{3}\times\sqrt{3}$\quad（2）$\sqrt{2}\times\sqrt{32}$\quad（3）$\sqrt{5}\times\sqrt{20}$
\bkpans{（1）\nobreak 3\quad（2）\nobreak $\sqrt{64}=8$\quad（3）\nobreak $\sqrt{100}=10$}

\bkproblem \R{長}{なが}さ\qty{50}{cm}の\R{振}{ふ}り\R{子}{こ}の\R{周期}{しゅう|き}は\R{約}{やく}1.4\R{秒}{びょう}です。\R{長}{なが}さ\qty{2}{m}の\R{振}{ふ}り\R{子}{こ}の\R{周期}{しゅう|き}は、\R{約}{やく}\R{何秒}{なん|びょう}ですか。
\bkpans{\qty{2}{m}＝\qty{200}{cm}で、$200\div50=4$\R{倍}{ばい}の\R{長}{なが}さ。\R{周期}{しゅう|き}は $\sqrt{4}=2$\R{倍}{ばい}で、$1.4\times2=2.8$。\R{約}{やく}2.8\R{秒}{びょう}}

\bkproblem \R{次}{つぎ}の\R{数}{かず}は、どの2つの\R{整数}{せい|すう}の\R{間}{あいだ}にありますか。\quad（1）$\sqrt{20}$\quad（2）$\sqrt{50}$
\bkpans{（1）\nobreak $4^2=16$、$5^2=25$ より4と5の\R{間}{あいだ}\quad（2）\nobreak $7^2=49$、$8^2=64$ より7と8の\R{間}{あいだ}}

\bkproblemgroup{チャレンジ}{2}

\bkproblem \R{正方形}{せい|ほう|けい}の\R{面積}{めん|せき}を9\R{倍}{ばい}にするには、1\R{辺}{ぺん}を\R{何倍}{なん|ばい}にすればよいですか。\R{面積}{めん|せき}を2\R{倍}{ばい}にするには、\R{約}{やく}\R{何倍}{なん|ばい}ですか。
\bkpans{$\sqrt{9}=3$\R{倍}{ばい}。2\R{倍}{ばい}にするには $\sqrt{2}$\R{倍}{ばい}で、\R{約}{やく}1.41\R{倍}{ばい}}

\bkproblem \R{長}{なが}さ\qty{25}{cm}で\R{周期}{しゅう|き}1.0\R{秒}{びょう}の\R{振}{ふ}り\R{子}{こ}があります。\R{周期}{しゅう|き}を2.0\R{秒}{びょう}にするには、\R{長}{なが}さを\R{何}{なん}cmにすればよいですか。
\bkpans{\R{周期}{しゅう|き}2\R{倍}{ばい}には、\R{長}{なが}さ $2^2=4$\R{倍}{ばい}。\qty{100}{cm}}

\bkproblem $\sqrt{2}\times\sqrt{2}$ と $\sqrt{4}$ は、どちらも2になります。$\sqrt{2}\times\sqrt{2}\times\sqrt{2}\times\sqrt{2}$ はいくつですか。
\bkpans{$(\sqrt{2}\times\sqrt{2})\times(\sqrt{2}\times\sqrt{2})=2\times2=4$}

\begin{column}{$\sqrt{\phantom{2}}$ の\R{値}{あたい}のおぼえ\R{方}{かた}}
\R{昔}{むかし}から、ごろ\R{合}{あ}わせでおぼえる\R{方法}{ほう|ほう}があります。\R{電卓}{でん|たく}で\R{計算}{けい|さん}して、\R{合}{あ}っているか\R{確}{たし}かめよう。
\begin{bkcenter}
\begin{bktable}{colspec={Q[l,m,wd=40mm] Q[l,m,wd=56mm]}, rowsep=1.5pt, colsep=4pt, row{1}={bg=bkNavy!8}}
\R{値}{あたい} & ごろ\R{合}{あ}わせ \\
$\sqrt{2}=1.41421356\cdots$ & ひとよひとよにひとみごろ \\
$\sqrt{3}=1.7320508\cdots$ & ひとなみにおごれや \\
$\sqrt{5}=2.2360679\cdots$ & ふじさんろくおうむなく \\
\end{bktable}
\end{bkcenter}
\end{column}

\begin{furikaeri}
\begin{checklist}
\item \R{面積}{めん|せき}から1\R{辺}{ぺん}を、2\R{乗}{じょう}ではさんで\R{求}{もと}められる。\hfill{\footnotesize\sffamily\color{bkTheme}→\pageref{bk:2-8:hasami}ページ}
\item $\sqrt{\phantom{2}}$ の\R{意味}{い|み}がわかり、\R{平方根}{へい|ほう|こん}を\R{求}{もと}められる。\hfill{\footnotesize\sffamily\color{bkTheme}→\pageref{bk:2-8:root}ページ}
\item \R{周期}{しゅう|き}が\R{長}{なが}さの $\sqrt{\phantom{2}}$ に\R{比例}{ひ|れい}することを\R{説明}{せつ|めい}できる。\hfill{\footnotesize\sffamily\color{bkTheme}→\pageref{bk:2-8:furiko}ページ}
\end{checklist}
\end{furikaeri}

\bkfukidashi{sora}{$\sqrt{2}$ は、どこまで\R{計算}{けい|さん}しても\R{終}{お}わらなかったね。\R{次}{つぎ}の2-9では、「\R{分数}{ぶん|すう}で\R{書}{か}けない\R{数}{かず}」のなぞにせまるよ。}
